% Manuscript support for results exported by the analysis code (generated copy of
% latex/setup.tex in the analysis repository; edit it there). Usage, in the preamble:
%   % Manuscript support for results exported by the analysis code (generated copy of
% latex/setup.tex in the analysis repository; edit it there). Usage, in the preamble:
%   % Manuscript support for results exported by the analysis code (generated copy of
% latex/setup.tex in the analysis repository; edit it there). Usage, in the preamble:
%   % Manuscript support for results exported by the analysis code (generated copy of
% latex/setup.tex in the analysis repository; edit it there). Usage, in the preamble:
%   \input{results/setup}
% Then:  \result[experiment_1, accuracy, flat]   \resulttable{name}   \autoresulttable{name}   \resultgraphic[width=\linewidth]{name}
\usepackage{pgfkeys}
\usepackage{graphicx}
\ExplSyntaxOn
\clist_new:N \l__results_keys_clist
\tl_new:N \l__results_path_tl

% \setresult{a/b/c}{value}: written by the analysis code (save_results()).
\NewDocumentCommand{\setresult}{mm}{ \pgfkeys{/results/#1/.initial={#2}} }

% \result[a, b, c]: the value stored under a/b/c. A hyphen in a key may be written for an underscore,
% e.g. \result[embedding-caption, global-lepi, text].
\NewDocumentCommand{\result}{r[]}
  {
    \clist_set:Nn \l__results_keys_clist {#1}
    \tl_set:Nx \l__results_path_tl
      { /results/ \clist_use:Nn \l__results_keys_clist { / } }
    % hyphens in the keys stand for underscores (avoids the editor's "_ must be inside math mode" warning)
    % (any `_` will do: only the character code counts in a control-sequence name)
    \tl_replace_all:Nnn \l__results_path_tl { - } { _ }
    \pgfkeysifdefined { \l__results_path_tl }
      { \exp_args:NV \pgfkeysvalueof \l__results_path_tl }
      {
        \msg_warning:nnx { results } { missing } { \l__results_path_tl }
        \textbf{??}
      }
  }
\msg_new:nnn { results } { missing } { Result~'#1'~is~not~defined. }
\ExplSyntaxOff

% \resulttable{name}: the tabular of results/tables/name.tex, to put in a table environment with its own
% placement, \caption and \label. \resultcaption{name}: the generated caption text of that table (to use or adapt).
\newcommand{\resulttable}[1]{\input{results/tables/#1}}
\newcommand{\resultcaption}[1]{\result[caption_#1, text]}
% \autoresulttable{name}[placement][label]: the whole table, i.e. a table environment with \centering, the
% tabular, the generated caption and \label{tab:name}. Optional: placement (default: none) and label
% (default: tab:name), e.g. \autoresulttable{calibration_verdicts}[htb][tab:calibration].
% Note: a "[" directly after the command is read as an option.
\NewDocumentCommand{\autoresulttable}{m o o}
  {
    \IfValueTF{#2}{\begin{table}[#2]}{\begin{table}}
      \centering
      \resulttable{#1}
      \caption{\resultcaption{#1}}
      \IfValueTF{#3}{\label{#3}}{\label{tab:#1}}
    \end{table}
  }
% \resultsection{name}: a section from results/sections/name.tex.
\newcommand{\resultsection}[1]{\input{results/sections/#1}}
% \resultgraphic[options]{name}: a figure from results/figures/name.pdf.
\newcommand{\resultgraphic}[2][]{\includegraphics[#1]{results/figures/#2}}

\IfFileExists{results/values/index.tex}{\input{results/values/index}}{}

% Then:  \result[experiment_1, accuracy, flat]   \resulttable{name}   \autoresulttable{name}   \resultgraphic[width=\linewidth]{name}
\usepackage{pgfkeys}
\usepackage{graphicx}
\ExplSyntaxOn
\clist_new:N \l__results_keys_clist
\tl_new:N \l__results_path_tl

% \setresult{a/b/c}{value}: written by the analysis code (save_results()).
\NewDocumentCommand{\setresult}{mm}{ \pgfkeys{/results/#1/.initial={#2}} }

% \result[a, b, c]: the value stored under a/b/c. A hyphen in a key may be written for an underscore,
% e.g. \result[embedding-caption, global-lepi, text].
\NewDocumentCommand{\result}{r[]}
  {
    \clist_set:Nn \l__results_keys_clist {#1}
    \tl_set:Nx \l__results_path_tl
      { /results/ \clist_use:Nn \l__results_keys_clist { / } }
    % hyphens in the keys stand for underscores (avoids the editor's "_ must be inside math mode" warning)
    % (any `_` will do: only the character code counts in a control-sequence name)
    \tl_replace_all:Nnn \l__results_path_tl { - } { _ }
    \pgfkeysifdefined { \l__results_path_tl }
      { \exp_args:NV \pgfkeysvalueof \l__results_path_tl }
      {
        \msg_warning:nnx { results } { missing } { \l__results_path_tl }
        \textbf{??}
      }
  }
\msg_new:nnn { results } { missing } { Result~'#1'~is~not~defined. }
\ExplSyntaxOff

% \resulttable{name}: the tabular of results/tables/name.tex, to put in a table environment with its own
% placement, \caption and \label. \resultcaption{name}: the generated caption text of that table (to use or adapt).
\newcommand{\resulttable}[1]{\input{results/tables/#1}}
\newcommand{\resultcaption}[1]{\result[caption_#1, text]}
% \autoresulttable{name}[placement][label]: the whole table, i.e. a table environment with \centering, the
% tabular, the generated caption and \label{tab:name}. Optional: placement (default: none) and label
% (default: tab:name), e.g. \autoresulttable{calibration_verdicts}[htb][tab:calibration].
% Note: a "[" directly after the command is read as an option.
\NewDocumentCommand{\autoresulttable}{m o o}
  {
    \IfValueTF{#2}{\begin{table}[#2]}{\begin{table}}
      \centering
      \resulttable{#1}
      \caption{\resultcaption{#1}}
      \IfValueTF{#3}{\label{#3}}{\label{tab:#1}}
    \end{table}
  }
% \resultsection{name}: a section from results/sections/name.tex.
\newcommand{\resultsection}[1]{\input{results/sections/#1}}
% \resultgraphic[options]{name}: a figure from results/figures/name.pdf.
\newcommand{\resultgraphic}[2][]{\includegraphics[#1]{results/figures/#2}}

\IfFileExists{results/values/index.tex}{\input{results/values/calibration_error}
\input{results/values/caption_appendix_calibration_global_lepi}
\input{results/values/caption_appendix_calibration_plantnet}
\input{results/values/caption_appendix_comparison_flemming_calibrated}
\input{results/values/caption_appendix_comparison_flemming_zero}
\input{results/values/caption_appendix_comparison_global_lepi_calibrated}
\input{results/values/caption_appendix_comparison_global_lepi_zero}
\input{results/values/caption_appendix_comparison_plantnet_calibrated}
\input{results/values/caption_appendix_comparison_plantnet_zero}
\input{results/values/caption_appendix_win_loss}
\input{results/values/caption_calibration_verdicts}
\input{results/values/caption_main_comparison_calibrated}
\input{results/values/caption_main_comparison_zero}
\input{results/values/embedding_caption}
\input{results/values/embedding_species_labels}
\input{results/values/embedding_taxa}
\input{results/values/hierarchy_contrast}
\input{results/values/main_comparison_win_loss}
\input{results/values/species_photo_contributors}
\input{results/values/taxonomic_alignment}
}{}

% Then:  \result[experiment_1, accuracy, flat]   \resulttable{name}   \autoresulttable{name}   \resultgraphic[width=\linewidth]{name}
\usepackage{pgfkeys}
\usepackage{graphicx}
\ExplSyntaxOn
\clist_new:N \l__results_keys_clist
\tl_new:N \l__results_path_tl

% \setresult{a/b/c}{value}: written by the analysis code (save_results()).
\NewDocumentCommand{\setresult}{mm}{ \pgfkeys{/results/#1/.initial={#2}} }

% \result[a, b, c]: the value stored under a/b/c. A hyphen in a key may be written for an underscore,
% e.g. \result[embedding-caption, global-lepi, text].
\NewDocumentCommand{\result}{r[]}
  {
    \clist_set:Nn \l__results_keys_clist {#1}
    \tl_set:Nx \l__results_path_tl
      { /results/ \clist_use:Nn \l__results_keys_clist { / } }
    % hyphens in the keys stand for underscores (avoids the editor's "_ must be inside math mode" warning)
    % (any `_` will do: only the character code counts in a control-sequence name)
    \tl_replace_all:Nnn \l__results_path_tl { - } { _ }
    \pgfkeysifdefined { \l__results_path_tl }
      { \exp_args:NV \pgfkeysvalueof \l__results_path_tl }
      {
        \msg_warning:nnx { results } { missing } { \l__results_path_tl }
        \textbf{??}
      }
  }
\msg_new:nnn { results } { missing } { Result~'#1'~is~not~defined. }
\ExplSyntaxOff

% \resulttable{name}: the tabular of results/tables/name.tex, to put in a table environment with its own
% placement, \caption and \label. \resultcaption{name}: the generated caption text of that table (to use or adapt).
\newcommand{\resulttable}[1]{\input{results/tables/#1}}
\newcommand{\resultcaption}[1]{\result[caption_#1, text]}
% \autoresulttable{name}[placement][label]: the whole table, i.e. a table environment with \centering, the
% tabular, the generated caption and \label{tab:name}. Optional: placement (default: none) and label
% (default: tab:name), e.g. \autoresulttable{calibration_verdicts}[htb][tab:calibration].
% Note: a "[" directly after the command is read as an option.
\NewDocumentCommand{\autoresulttable}{m o o}
  {
    \IfValueTF{#2}{\begin{table}[#2]}{\begin{table}}
      \centering
      \resulttable{#1}
      \caption{\resultcaption{#1}}
      \IfValueTF{#3}{\label{#3}}{\label{tab:#1}}
    \end{table}
  }
% \resultsection{name}: a section from results/sections/name.tex.
\newcommand{\resultsection}[1]{\input{results/sections/#1}}
% \resultgraphic[options]{name}: a figure from results/figures/name.pdf.
\newcommand{\resultgraphic}[2][]{\includegraphics[#1]{results/figures/#2}}

\IfFileExists{results/values/index.tex}{% Generated; do not edit.
\setresult{calibration_error/global_lepi/higher/flat/error}{4.28}
\setresult{calibration_error/global_lepi/higher/hierarchical/error}{1.63}
\setresult{calibration_error/global_lepi/higher/hierarchical/difference}{$-2.65$ ($-2.70$ to $-2.59$)}
\setresult{calibration_error/global_lepi/higher/hierarchical/verdict}{better}
\setresult{calibration_error/global_lepi/higher/hierarchical/backbone_verdicts}{3 better}
\setresult{calibration_error/global_lepi/higher/conditional/error}{2.28}
\setresult{calibration_error/global_lepi/higher/conditional/difference}{$-2.00$ ($-2.06$ to $-1.93$)}
\setresult{calibration_error/global_lepi/higher/conditional/verdict}{better}
\setresult{calibration_error/global_lepi/higher/conditional/backbone_verdicts}{3 better}
\setresult{calibration_error/global_lepi/higher/independent/error}{2.97}
\setresult{calibration_error/global_lepi/higher/independent/difference}{$-1.31$ ($-1.36$ to $-1.25$)}
\setresult{calibration_error/global_lepi/higher/independent/verdict}{better}
\setresult{calibration_error/global_lepi/higher/independent/backbone_verdicts}{3 better}
\setresult{calibration_error/global_lepi/species/flat/error}{2.39}
\setresult{calibration_error/global_lepi/species/hierarchical/error}{1.91}
\setresult{calibration_error/global_lepi/species/hierarchical/difference}{$-0.48$ ($-0.55$ to $-0.40$)}
\setresult{calibration_error/global_lepi/species/hierarchical/verdict}{better}
\setresult{calibration_error/global_lepi/species/hierarchical/backbone_verdicts}{2 better and 1 worse}
\setresult{calibration_error/global_lepi/species/conditional/error}{5.09}
\setresult{calibration_error/global_lepi/species/conditional/difference}{$2.71$ ($2.63$ to $2.79$)}
\setresult{calibration_error/global_lepi/species/conditional/verdict}{worse}
\setresult{calibration_error/global_lepi/species/conditional/backbone_verdicts}{1 better and 2 worse}
\setresult{calibration_error/global_lepi/species/independent/error}{2.45}
\setresult{calibration_error/global_lepi/species/independent/difference}{$0.06$ ($-0.01$ to $0.14$)}
\setresult{calibration_error/global_lepi/species/independent/verdict}{inconclusive}
\setresult{calibration_error/global_lepi/species/independent/backbone_verdicts}{1 better, 1 worse and 1 inconclusive}
\setresult{calibration_error/global_lepi/backbones}{3}
\setresult{calibration_error/plantnet/higher/flat/error}{5.24}
\setresult{calibration_error/plantnet/higher/hierarchical/error}{2.08}
\setresult{calibration_error/plantnet/higher/hierarchical/difference}{$-3.16$ ($-3.62$ to $-2.70$)}
\setresult{calibration_error/plantnet/higher/hierarchical/verdict}{better}
\setresult{calibration_error/plantnet/higher/hierarchical/backbone_verdicts}{2 better and 1 inconclusive}
\setresult{calibration_error/plantnet/higher/conditional/error}{3.09}
\setresult{calibration_error/plantnet/higher/conditional/difference}{$-2.15$ ($-2.57$ to $-1.72$)}
\setresult{calibration_error/plantnet/higher/conditional/verdict}{better}
\setresult{calibration_error/plantnet/higher/conditional/backbone_verdicts}{2 better and 1 worse}
\setresult{calibration_error/plantnet/higher/independent/error}{2.61}
\setresult{calibration_error/plantnet/higher/independent/difference}{$-2.63$ ($-3.06$ to $-2.20$)}
\setresult{calibration_error/plantnet/higher/independent/verdict}{better}
\setresult{calibration_error/plantnet/higher/independent/backbone_verdicts}{2 better and 1 inconclusive}
\setresult{calibration_error/plantnet/species/flat/error}{3.99}
\setresult{calibration_error/plantnet/species/hierarchical/error}{3.36}
\setresult{calibration_error/plantnet/species/hierarchical/difference}{$-0.63$ ($-1.03$ to $-0.23$)}
\setresult{calibration_error/plantnet/species/hierarchical/verdict}{better}
\setresult{calibration_error/plantnet/species/hierarchical/backbone_verdicts}{2 better and 1 inconclusive}
\setresult{calibration_error/plantnet/species/conditional/error}{11.85}
\setresult{calibration_error/plantnet/species/conditional/difference}{$7.86$ ($7.44$ to $8.28$)}
\setresult{calibration_error/plantnet/species/conditional/verdict}{worse}
\setresult{calibration_error/plantnet/species/conditional/backbone_verdicts}{3 worse}
\setresult{calibration_error/plantnet/species/independent/error}{5.33}
\setresult{calibration_error/plantnet/species/independent/difference}{$1.34$ ($0.94$ to $1.74$)}
\setresult{calibration_error/plantnet/species/independent/verdict}{worse}
\setresult{calibration_error/plantnet/species/independent/backbone_verdicts}{2 worse and 1 inconclusive}
\setresult{calibration_error/plantnet/backbones}{3}
\setresult{calibration_error/overall/hierarchical}{9 better, 1 worse and 2 inconclusive}
\setresult{calibration_error/overall/conditional}{6 better and 6 worse}
\setresult{calibration_error/overall/independent}{6 better, 3 worse and 3 inconclusive}

% Generated; do not edit.
\setresult{caption_appendix_calibration_global_lepi/text}{Calibration error (\%) of the heads in the Global Lepidoptera dataset, per rank and backbone. Calibration error is the square root of the debiased mean squared accuracy--confidence gap across confidence bins. Better/worse: the pointwise 95\% confidence interval for the paired difference in squared calibration error lies entirely below/above 0; inconclusive: it includes 0. Bold: better calibrated than the flat head; italics: worse calibrated; plain entries (not including the flat column): inconclusive. Nonpositive estimates of the squared calibration error are displayed as zero; this does not establish exact calibration.}

% Generated; do not edit.
\setresult{caption_appendix_calibration_plantnet/text}{Calibration error (\%) of the heads in the Pl@ntNet300K dataset, per rank and backbone. Calibration error is the square root of the debiased mean squared accuracy--confidence gap across confidence bins. Better/worse: the pointwise 95\% confidence interval for the paired difference in squared calibration error lies entirely below/above 0; inconclusive: it includes 0. Bold: better calibrated than the flat head; italics: worse calibrated; plain entries (not including the flat column): inconclusive. Nonpositive estimates of the squared calibration error are displayed as zero; this does not establish exact calibration.}

% Generated; do not edit.
\setresult{caption_appendix_comparison_flemming_calibrated/text}{Test metrics for Expert labeled AMI (calibrated thresholds), in percent: estimate on the reporting partition (90\% of the test set) $\pm$ standard deviation over 1000 bootstrap resamples. Bold: the head is a plausible winner, i.e. its share of resamples in which it is best plausibly reaches chance ($1/K$ for $K$ heads; posterior probability above 5\%). Underlined: clear winner, i.e. every other head is below chance with posterior probability above 95\%. The probabilities concern the resampling shares given the fixed test data and predictions. Wins are decided on metrics rounded to 0.1 percentage points. Compared within each backbone and rank.}

% Generated; do not edit.
\setresult{caption_appendix_comparison_flemming_zero/text}{Test metrics for Expert labeled AMI (zero thresholds), in percent: estimate on the reporting partition (90\% of the test set) $\pm$ standard deviation over 1000 bootstrap resamples. Bold: the head is a plausible winner, i.e. its share of resamples in which it is best plausibly reaches chance ($1/K$ for $K$ heads; posterior probability above 5\%). Underlined: clear winner, i.e. every other head is below chance with posterior probability above 95\%. The probabilities concern the resampling shares given the fixed test data and predictions. Wins are decided on metrics rounded to 0.1 percentage points. Compared within each backbone and rank.}

% Generated; do not edit.
\setresult{caption_appendix_comparison_global_lepi_calibrated/text}{Test metrics for Global Lepidoptera (calibrated thresholds), in percent: estimate on the reporting partition (90\% of the test set) $\pm$ standard deviation over 200 bootstrap resamples. Bold: the head is a plausible winner, i.e. its share of resamples in which it is best plausibly reaches chance ($1/K$ for $K$ heads; posterior probability above 5\%). Underlined: clear winner, i.e. every other head is below chance with posterior probability above 95\%. The probabilities concern the resampling shares given the fixed test data and predictions. Wins are decided on metrics rounded to 0.1 percentage points. Compared within each backbone and rank.}

% Generated; do not edit.
\setresult{caption_appendix_comparison_global_lepi_zero/text}{Test metrics for Global Lepidoptera (zero thresholds), in percent: estimate on the reporting partition (90\% of the test set) $\pm$ standard deviation over 200 bootstrap resamples. Bold: the head is a plausible winner, i.e. its share of resamples in which it is best plausibly reaches chance ($1/K$ for $K$ heads; posterior probability above 5\%). Underlined: clear winner, i.e. every other head is below chance with posterior probability above 95\%. The probabilities concern the resampling shares given the fixed test data and predictions. Wins are decided on metrics rounded to 0.1 percentage points. Compared within each backbone and rank.}

% Generated; do not edit.
\setresult{caption_appendix_comparison_plantnet_calibrated/text}{Test metrics for Pl@ntNet300K (calibrated thresholds), in percent: estimate on the reporting partition (90\% of the test set) $\pm$ standard deviation over 1000 bootstrap resamples. Bold: the head is a plausible winner, i.e. its share of resamples in which it is best plausibly reaches chance ($1/K$ for $K$ heads; posterior probability above 5\%). Underlined: clear winner, i.e. every other head is below chance with posterior probability above 95\%. The probabilities concern the resampling shares given the fixed test data and predictions. Wins are decided on metrics rounded to 0.1 percentage points. Compared within each backbone and rank.}

% Generated; do not edit.
\setresult{caption_appendix_comparison_plantnet_zero/text}{Test metrics for Pl@ntNet300K (zero thresholds), in percent: estimate on the reporting partition (90\% of the test set) $\pm$ standard deviation over 1000 bootstrap resamples. Bold: the head is a plausible winner, i.e. its share of resamples in which it is best plausibly reaches chance ($1/K$ for $K$ heads; posterior probability above 5\%). Underlined: clear winner, i.e. every other head is below chance with posterior probability above 95\%. The probabilities concern the resampling shares given the fixed test data and predictions. Wins are decided on metrics rounded to 0.1 percentage points. Compared within each backbone and rank.}

% Generated; do not edit.
\setresult{caption_appendix_win_loss/text}{Share (\%) of comparisons in which each head is a winner or a loser, and a clear winner or loser (see Methods), at the species, genus and family rank. A comparison is one dataset, backbone, setting and metric (n = 108 per rank, excluding coverage). Mean: average over the three ranks.}

% Generated; do not edit.
\setresult{caption_calibration_verdicts/text}{Calibration of the heads compared with the flat head: the number of cases (3 backbones in each of 2 datasets) in which a head is better or worse calibrated than the flat head, at the species rank and pooled over the higher ranks. Calibration error is the square root of the debiased mean squared accuracy--confidence gap across confidence bins. Better/worse: the pointwise 95\% confidence interval for the paired difference in squared calibration error lies entirely below/above 0; inconclusive: it includes 0.}

% Generated; do not edit.
\setresult{caption_main_comparison_calibrated/text}{Test metrics for EfficientNetV2-Small (calibrated thresholds), in percent: estimate on the reporting partition (90\% of the test set) $\pm$ standard deviation over 1000 (Pl@ntNet300K and Expert labeled AMI) or 200 (Global Lepidoptera) bootstrap resamples. Bold: the head is a plausible winner, i.e. its share of resamples in which it is best plausibly reaches chance ($1/K$ for $K$ heads; posterior probability above 5\%). Underlined: clear winner, i.e. every other head is below chance with posterior probability above 95\%. The probabilities concern the resampling shares given the fixed test data and predictions. Wins are decided on metrics rounded to 0.1 percentage points. Compared within each dataset and rank.}

% Generated; do not edit.
\setresult{caption_main_comparison_zero/text}{Test metrics for EfficientNetV2-Small (zero thresholds), in percent: estimate on the reporting partition (90\% of the test set) $\pm$ standard deviation over 1000 (Pl@ntNet300K and Expert labeled AMI) or 200 (Global Lepidoptera) bootstrap resamples. Bold: the head is a plausible winner, i.e. its share of resamples in which it is best plausibly reaches chance ($1/K$ for $K$ heads; posterior probability above 5\%). Underlined: clear winner, i.e. every other head is below chance with posterior probability above 95\%. The probabilities concern the resampling shares given the fixed test data and predictions. Wins are decided on metrics rounded to 0.1 percentage points. Compared within each dataset and rank.}

% Generated; do not edit.
\setresult{embedding_caption/plantnet/species}{34}
\setresult{embedding_caption/plantnet/rare}{11}
\setresult{embedding_caption/plantnet/without_photo}{1}
\setresult{embedding_caption/plantnet/cap}{30}
\setresult{embedding_caption/plantnet/min_images}{4}
\setresult{embedding_caption/plantnet/text}{Embedding structure of the 34 species of the family \textit{Ranunculaceae} in the Pl@ntNet300K dataset: up to 30 test images per species and the species' class means, at the angle to the family's spherical median (rings every 15 degrees). The radial axis of each panel runs from 50 to 75 degrees (Flat), 45 to 75 degrees (Hierarchical), 45 to 70 degrees (Independent) and 5 to 45 degrees (Conditional); images beyond the outer limit are placed on it. Hollow markers (faded in the photo figures) show the 11 species with fewer than 4 images. No photo is available for 1 species.}
\setresult{embedding_caption/global_lepi/species}{68}
\setresult{embedding_caption/global_lepi/rare}{0}
\setresult{embedding_caption/global_lepi/without_photo}{2}
\setresult{embedding_caption/global_lepi/cap}{30}
\setresult{embedding_caption/global_lepi/min_images}{4}
\setresult{embedding_caption/global_lepi/text}{Embedding structure of the 68 species of the family \textit{Pterophoridae} in the Global Lepidoptera dataset: up to 30 test images per species and the species' class means, at the angle to the family's spherical median (rings every 15 degrees). The radial axis of each panel runs from 30 to 90 degrees (Flat), 30 to 90 degrees (Hierarchical), 35 to 80 degrees (Independent) and 30 to 75 degrees (Conditional); images beyond the outer limit are placed on it. No photo is available for 2 species.}

% Generated; do not edit.
\setresult{embedding_species_labels/plantnet}{Species labels in the Pl@ntNet300K photo figures, with the full names and GBIF taxon keys of the species: \textit{A. aestivalis} = \textit{Adonis aestivalis} (5371707); \textit{A. annua} = \textit{Adonis annua} (5371709); \textit{A. flammea} = \textit{Adonis flammea} (8246651); \textit{A. microcarpa} = \textit{Adonis microcarpa} (5616463); \textit{A. pyrenaica} = \textit{Adonis pyrenaica} (8158904); \textit{A. vernalis} = \textit{Adonis vernalis} (5371711); \textit{A. canadense} = \textit{Anemonastrum canadense} (9419395); \textit{A. narcissiflorum} = \textit{Anemonastrum narcissiflorum} (8870733); \textit{A. apennina} = \textit{Anemone apennina} (3927746); \textit{A. baldensis} = \textit{Anemone baldensis} (3927618); \textit{A. blanda} = \textit{Anemone blanda} (3033242); \textit{A. coronaria} = \textit{Anemone coronaria} (3033261); \textit{A. hortensis} = \textit{Anemone hortensis} (3925918); \textit{A. multifida} = \textit{Anemone multifida} (3033252); \textit{A. nemorosa} = \textit{Anemone nemorosa} (3033263); \textit{A. palmata} = \textit{Anemone palmata} (3924291); \textit{A. ranunculoides} = \textit{Anemone ranunculoides} (3033266); \textit{A. sylvestris} = \textit{Anemone sylvestris} (3922888); \textit{A. trifolia} = \textit{Anemone trifolia} (3033231); \textit{A. virginiana} = \textit{Anemone virginiana} (3033255); \textit{E. hupehensis} = \textit{Eriocapitella hupehensis} (9677963); \textit{E. hybrida} = \textit{Eriocapitella hybrida} (9673106); \textit{E. vitifolia} = \textit{Eriocapitella vitifolia} (3925297); \textit{H. nobilis} = \textit{Hepatica nobilis} (5371699); \textit{M. minimus} = \textit{Myosurus minimus} (3033605); \textit{P. alpina} = \textit{Pulsatilla alpina} (7278071); \textit{P. halleri} = \textit{Pulsatilla halleri} (3929423); \textit{P. montana} = \textit{Pulsatilla montana} (7277925); \textit{P. patens} = \textit{Pulsatilla patens} (3033623); \textit{P. pratensis} = \textit{Pulsatilla pratensis} (7277878); \textit{P. rubra} = \textit{Pulsatilla rubra} (7361164); \textit{P. vernalis} = \textit{Pulsatilla vernalis} (3928615); \textit{P. vulgaris} = \textit{Pulsatilla vulgaris} (3033621).}
\setresult{embedding_species_labels/global_lepi}{Species labels in the Global Lepidoptera photo figures, with the full names and GBIF taxon keys of the species: \textit{A. ambrosiae} = \textit{Adaina ambrosiae} (1858706); \textit{A. microdactyla} = \textit{Adaina microdactyla} (1858698); \textit{A. montanus} = \textit{Adaina montanus} (9689122); \textit{A. acanthadactyla} = \textit{Amblyptilia acanthadactyla} (1859285); \textit{A. falcatalis} = \textit{Amblyptilia falcatalis} (9590272); \textit{A. pica} = \textit{Amblyptilia pica} (1859287); \textit{A. punctidactyla} = \textit{Amblyptilia punctidactyla} (1859305); \textit{A. repletalis} = \textit{Amblyptilia repletalis} (9801765); \textit{A. marmarodactyla} = \textit{Anstenoptilia marmarodactyla} (1858615); \textit{C. rhododactyla} = \textit{Cnaemidophorus rhododactyla} (1858637); \textit{C. aglaodesma} = \textit{Cosmoclostis aglaodesma} (1858973); \textit{C. distans} = \textit{Crombrugghia distans} (1860063); \textit{C. tristis} = \textit{Crombrugghia tristis} (1860076); \textit{D. californicus} = \textit{Dejongia californicus} (1858659); \textit{D. lobidactylus} = \textit{Dejongia lobidactylus} (1858660); \textit{D. fasciola} = \textit{Diacrotricha fasciola} (1859453); \textit{E. monodactyla} = \textit{Emmelina monodactyla} (1859046); \textit{E. pumilio} = \textit{Exelastis pumilio} (6133397); \textit{G. didactyla} = \textit{Geina didactyla} (5123042); \textit{G. periscelidactylus} = \textit{Geina periscelidactylus} (5123046); \textit{G. ochrodactyla} = \textit{Gillmeria ochrodactyla} (4536483); \textit{G. pallidactyla} = \textit{Gillmeria pallidactyla} (4536472); \textit{H. balanotes} = \textit{Hellinsia balanotes} (9675743); \textit{H. chamelai} = \textit{Hellinsia chamelai} (8790054); \textit{H. elliottii} = \textit{Hellinsia elliottii} (9746673); \textit{H. homodactylus} = \textit{Hellinsia homodactylus} (9748079); \textit{H. inquinatus} = \textit{Hellinsia inquinatus} (9686431); \textit{H. osteodactylus} = \textit{Hellinsia osteodactylus} (1859480); \textit{H. tephradactyla} = \textit{Hellinsia tephradactyla} (1859488); \textit{L. pusillidactylus} = \textit{Lantanophaga pusillidactylus} (1858674); \textit{L. albistriolatus} = \textit{Lioptilodes albistriolatus} (9760511); \textit{M. lunaedactyla} = \textit{Marasmarcha lunaedactyla} (1859534); \textit{M. leucodactylus} = \textit{Megalorhipida leucodactylus} (1858930); \textit{M. baliodactylus} = \textit{Merrifieldia baliodactylus} (7642240); \textit{M. leucodactyla} = \textit{Merrifieldia leucodactyla} (8200050); \textit{M. indentatus} = \textit{Michaelophorus indentatus} (9700261); \textit{N. cinctipedalis} = \textit{Nippoptilia cinctipedalis} (9749806); \textit{O. eupatorii} = \textit{Oidaematophorus eupatorii} (9752950); \textit{O. lithodactyla} = \textit{Oidaematophorus lithodactyla} (1859952); \textit{O. pilosellae} = \textit{Oxyptilus pilosellae} (1858873); \textit{P. chlorosata} = \textit{Petrophora chlorosata} (1973467); \textit{P. convergata} = \textit{Petrophora convergata} (4525132); \textit{P. divisata} = \textit{Petrophora divisata} (1973458); \textit{P. narbonea} = \textit{Petrophora narbonea} (12046692); \textit{P. subaequaria} = \textit{Petrophora subaequaria} (1973466); \textit{P. calodactyla} = \textit{Platyptilia calodactyla} (1859100); \textit{P. carduidactylus} = \textit{Platyptilia carduidactylus} (1859161); \textit{P. celidotus} = \textit{Platyptilia celidotus} (9768621); \textit{P. gonodactyla} = \textit{Platyptilia gonodactyla} (1859177); \textit{P. isodactylus} = \textit{Platyptilia isodactylus} (4536292); \textit{P. belfragei} = \textit{Pselnophorus belfragei} (9777614); \textit{P. furcatalis} = \textit{Pterophorus furcatalis} (5123118); \textit{P. innotatalis} = \textit{Pterophorus innotatalis} (5123137); \textit{P. lacteipennis} = \textit{Pterophorus lacteipennis} (5123180); \textit{P. monospilalis} = \textit{Pterophorus monospilalis} (5123276); \textit{P. pentadactyla} = \textit{Pterophorus pentadactyla} (5123221); \textit{S. alternaria} = \textit{Singularia alternaria} (9645309); \textit{S. emissalis} = \textit{Sinpunctiptilia emissalis} (9759340); \textit{S. anisodactylus} = \textit{Sphenarches anisodactylus} (1859932); \textit{S. ontario} = \textit{Sphenarches ontario} (1859929); \textit{S. xerodes} = \textit{Stangeia xerodes} (1859061); \textit{S. bipunctidactyla} = \textit{Stenoptilia bipunctidactyla} (1858778); \textit{S. pterodactyla} = \textit{Stenoptilia pterodactyla} (1858775); \textit{S. zophodactylus} = \textit{Stenoptilia zophodactylus} (1858762); \textit{S. taprobanes} = \textit{Stenoptilodes taprobanes} (1860026); \textit{W. spilodactylus} = \textit{Wheeleria spilodactylus} (6097269).}

% Generated; do not edit.
\setresult{embedding_taxa/plantnet}{Taxon names in the Pl@ntNet300K figures with their GBIF taxon keys: the family \textit{Ranunculaceae} (2410) and the genera \textit{Adonis} (3033821), \textit{Anemonastrum} (7277728), \textit{Anemone} (3016227), \textit{Eriocapitella} (9777080), \textit{Hepatica} (3033655), \textit{Myosurus} (3033602) and \textit{Pulsatilla} (3033620).}
\setresult{embedding_taxa/global_lepi}{Taxon names in the Global Lepidoptera figures with their GBIF taxon keys: the family \textit{Pterophoridae} (8863) and the genera \textit{Adaina} (1858684), \textit{Amblyptilia} (1859284), \textit{Dejongia} (1858658), \textit{Geina} (9077410), \textit{Gillmeria} (4406145), \textit{Hellinsia} (1859454), \textit{Oidaematophorus} (1859936), \textit{Platyptilia} (1859082), \textit{Pterophorus} (8331787), \textit{Sphenarches} (1859925) and \textit{Stenoptilia} (1858722). The genera grouped as Other are \textit{Anstenoptilia} (1858614), \textit{Cnaemidophorus} (1858636), \textit{Cosmoclostis} (1858971), \textit{Crombrugghia} (1860060), \textit{Diacrotricha} (1859451), \textit{Emmelina} (1859036), \textit{Exelastis} (1859547), \textit{Hexadactilia} (1858716), \textit{Lantanophaga} (1858672), \textit{Lioptilodes} (1859311), \textit{Marasmarcha} (1859503), \textit{Megalorhipida} (1858921), \textit{Merrifieldia} (7641205), \textit{Michaelophorus} (4708885), \textit{Nippoptilia} (1859540), \textit{Oxyptilus} (1858866), \textit{Pselnophorus} (1858640), \textit{Singularia} (1859070), \textit{Sinpunctiptilia} (9260167), \textit{Stangeia} (1859059), \textit{Stenoptilodes} (1860010) and \textit{Wheeleria} (8744407).}

% Generated; do not edit.
\setresult{hierarchy_contrast/images}{31088}
\setresult{hierarchy_contrast/effect/architecture}{-2.35}
\setresult{hierarchy_contrast/effect/seed42}{-2.01}
\setresult{hierarchy_contrast/effect/seed43}{-2.35}
\setresult{hierarchy_contrast/effect/seed44}{-2.34}
\setresult{hierarchy_contrast/difference/estimate}{-0.12}
\setresult{hierarchy_contrast/difference/lower}{-0.54}
\setresult{hierarchy_contrast/difference/upper}{0.30}

% Generated; do not edit.
\setresult{main_comparison_win_loss/conditional/family/wins}{19.44}
\setresult{main_comparison_win_loss/conditional/family/clear_winner}{4.63}
\setresult{main_comparison_win_loss/conditional/family/winner}{14.81}
\setresult{main_comparison_win_loss/conditional/family/losses}{43.52}
\setresult{main_comparison_win_loss/conditional/family/clear_loser}{23.15}
\setresult{main_comparison_win_loss/conditional/family/loser}{20.37}
\setresult{main_comparison_win_loss/conditional/genus/wins}{21.30}
\setresult{main_comparison_win_loss/conditional/genus/clear_winner}{7.41}
\setresult{main_comparison_win_loss/conditional/genus/winner}{13.89}
\setresult{main_comparison_win_loss/conditional/genus/losses}{53.70}
\setresult{main_comparison_win_loss/conditional/genus/clear_loser}{40.74}
\setresult{main_comparison_win_loss/conditional/genus/loser}{12.96}
\setresult{main_comparison_win_loss/conditional/mean/wins}{19.14}
\setresult{main_comparison_win_loss/conditional/mean/clear_winner}{5.86}
\setresult{main_comparison_win_loss/conditional/mean/winner}{13.27}
\setresult{main_comparison_win_loss/conditional/mean/losses}{55.56}
\setresult{main_comparison_win_loss/conditional/mean/clear_loser}{41.98}
\setresult{main_comparison_win_loss/conditional/mean/loser}{13.58}
\setresult{main_comparison_win_loss/conditional/species/wins}{16.67}
\setresult{main_comparison_win_loss/conditional/species/clear_winner}{5.56}
\setresult{main_comparison_win_loss/conditional/species/winner}{11.11}
\setresult{main_comparison_win_loss/conditional/species/losses}{69.44}
\setresult{main_comparison_win_loss/conditional/species/clear_loser}{62.04}
\setresult{main_comparison_win_loss/conditional/species/loser}{7.41}
\setresult{main_comparison_win_loss/flat/family/wins}{22.22}
\setresult{main_comparison_win_loss/flat/family/clear_winner}{9.26}
\setresult{main_comparison_win_loss/flat/family/winner}{12.96}
\setresult{main_comparison_win_loss/flat/family/losses}{48.15}
\setresult{main_comparison_win_loss/flat/family/clear_loser}{23.15}
\setresult{main_comparison_win_loss/flat/family/loser}{25.00}
\setresult{main_comparison_win_loss/flat/genus/wins}{16.67}
\setresult{main_comparison_win_loss/flat/genus/clear_winner}{5.56}
\setresult{main_comparison_win_loss/flat/genus/winner}{11.11}
\setresult{main_comparison_win_loss/flat/genus/losses}{41.67}
\setresult{main_comparison_win_loss/flat/genus/clear_loser}{29.63}
\setresult{main_comparison_win_loss/flat/genus/loser}{12.04}
\setresult{main_comparison_win_loss/flat/mean/wins}{28.40}
\setresult{main_comparison_win_loss/flat/mean/clear_winner}{9.88}
\setresult{main_comparison_win_loss/flat/mean/winner}{18.52}
\setresult{main_comparison_win_loss/flat/mean/losses}{39.51}
\setresult{main_comparison_win_loss/flat/mean/clear_loser}{23.77}
\setresult{main_comparison_win_loss/flat/mean/loser}{15.74}
\setresult{main_comparison_win_loss/flat/species/wins}{46.30}
\setresult{main_comparison_win_loss/flat/species/clear_winner}{14.81}
\setresult{main_comparison_win_loss/flat/species/winner}{31.48}
\setresult{main_comparison_win_loss/flat/species/losses}{28.70}
\setresult{main_comparison_win_loss/flat/species/clear_loser}{18.52}
\setresult{main_comparison_win_loss/flat/species/loser}{10.19}
\setresult{main_comparison_win_loss/hierarchical/family/wins}{64.81}
\setresult{main_comparison_win_loss/hierarchical/family/clear_winner}{26.85}
\setresult{main_comparison_win_loss/hierarchical/family/winner}{37.96}
\setresult{main_comparison_win_loss/hierarchical/family/losses}{1.85}
\setresult{main_comparison_win_loss/hierarchical/family/clear_loser}{1.85}
\setresult{main_comparison_win_loss/hierarchical/family/loser}{0.00}
\setresult{main_comparison_win_loss/hierarchical/genus/wins}{58.33}
\setresult{main_comparison_win_loss/hierarchical/genus/clear_winner}{29.63}
\setresult{main_comparison_win_loss/hierarchical/genus/winner}{28.70}
\setresult{main_comparison_win_loss/hierarchical/genus/losses}{6.48}
\setresult{main_comparison_win_loss/hierarchical/genus/clear_loser}{2.78}
\setresult{main_comparison_win_loss/hierarchical/genus/loser}{3.70}
\setresult{main_comparison_win_loss/hierarchical/mean/wins}{54.63}
\setresult{main_comparison_win_loss/hierarchical/mean/clear_winner}{22.84}
\setresult{main_comparison_win_loss/hierarchical/mean/winner}{31.79}
\setresult{main_comparison_win_loss/hierarchical/mean/losses}{8.02}
\setresult{main_comparison_win_loss/hierarchical/mean/clear_loser}{3.09}
\setresult{main_comparison_win_loss/hierarchical/mean/loser}{4.94}
\setresult{main_comparison_win_loss/hierarchical/species/wins}{40.74}
\setresult{main_comparison_win_loss/hierarchical/species/clear_winner}{12.04}
\setresult{main_comparison_win_loss/hierarchical/species/winner}{28.70}
\setresult{main_comparison_win_loss/hierarchical/species/losses}{15.74}
\setresult{main_comparison_win_loss/hierarchical/species/clear_loser}{4.63}
\setresult{main_comparison_win_loss/hierarchical/species/loser}{11.11}
\setresult{main_comparison_win_loss/independent/family/wins}{37.96}
\setresult{main_comparison_win_loss/independent/family/clear_winner}{16.67}
\setresult{main_comparison_win_loss/independent/family/winner}{21.30}
\setresult{main_comparison_win_loss/independent/family/losses}{35.19}
\setresult{main_comparison_win_loss/independent/family/clear_loser}{25.00}
\setresult{main_comparison_win_loss/independent/family/loser}{10.19}
\setresult{main_comparison_win_loss/independent/genus/wins}{48.15}
\setresult{main_comparison_win_loss/independent/genus/clear_winner}{16.67}
\setresult{main_comparison_win_loss/independent/genus/winner}{31.48}
\setresult{main_comparison_win_loss/independent/genus/losses}{14.81}
\setresult{main_comparison_win_loss/independent/genus/clear_loser}{10.19}
\setresult{main_comparison_win_loss/independent/genus/loser}{4.63}
\setresult{main_comparison_win_loss/independent/mean/wins}{49.07}
\setresult{main_comparison_win_loss/independent/mean/clear_winner}{14.81}
\setresult{main_comparison_win_loss/independent/mean/winner}{34.26}
\setresult{main_comparison_win_loss/independent/mean/losses}{18.21}
\setresult{main_comparison_win_loss/independent/mean/clear_loser}{11.73}
\setresult{main_comparison_win_loss/independent/mean/loser}{6.48}
\setresult{main_comparison_win_loss/independent/species/wins}{61.11}
\setresult{main_comparison_win_loss/independent/species/clear_winner}{11.11}
\setresult{main_comparison_win_loss/independent/species/winner}{50.00}
\setresult{main_comparison_win_loss/independent/species/losses}{4.63}
\setresult{main_comparison_win_loss/independent/species/clear_loser}{0.00}
\setresult{main_comparison_win_loss/independent/species/loser}{4.63}

% Generated; do not edit.
\setresult{species_photo_contributors/plantnet}{Adam Golenia, Alena Shindyapkina, Antoine MARNAT, Christian Berg, Daniel Ellehammer Larsen, David Delon, Dennis Lausch, Denys Vynokurov, Flavien Saboureau, Galyna Mykytynets, Jason Grant, Jonathan, Katrin Simon, Kristof Zyskowski, Ljaž, Marco Mussita, Matt Berger, Oleksii Vasyliuk, Roger Butterfield, Roman A. Lagazyuk (Роман А. Лагазюк), Shaun Coleman, Tatsutomo Chin, Vicente Lopes, Yves Bas, carnifex, dcourtine, joan\_carles, mister\_bumble, mykhas, onotole, phiro, radoslawpuchalka and terrikoontz}
\setresult{species_photo_contributors/global_lepi}{Aalbert Rebergen, Alexis Tinker-Tsavalas, Amber M. King, Andrew Meeds, Andrzej Chruślak, Annika Lindqvist, Cathy Powers, Dan MacNeal, Dan Riley, David George, Diego González Dopico, Donald Hobern, Doreen, Evgenij Samarin, Francisco Farriols Sarabia, Fyodor Pudovikov, Jason Grant, Joey Bom, Johannes S., Kai Squires, Keith Wauthy, Ken-ichi Ueda, Laura Gaudette, Louis Imbeau, Michael Berardozzi, Nigel Voaden, Peter Crowcroft, RAP, Randy A Nonenmacher, Romilson Silva Lopes Junior, Sam Kieschnick, Saryu Mae 前 朝琉, Seig, Shannon Stoner, Stuart Fraser, Sylvain Montagner, Thomas Irvine, Thomas Koffel, Tommy Andriollo, adam1420, amzamz, bioshots\_jm, carey\_knox\_southern\_scales, kike Vergara, loarie, naturewatchwidow, pmrsco, portioid, rich\_fuller, sgalick, steve\_mcwilliam, Вячеслав Юсупов and Поляков Александр}

% Generated; do not edit.
\setresult{taxonomic_alignment/methods}{average linkage, complete linkage, PAM and HDBSCAN}
\setresult{taxonomic_alignment/subsamples}{1000}
\setresult{taxonomic_alignment/species_per_subsample}{512}
\setresult{taxonomic_alignment/plantnet/genus/hierarchical/median}{0.15}
\setresult{taxonomic_alignment/plantnet/genus/hierarchical/lower}{0.02}
\setresult{taxonomic_alignment/plantnet/genus/hierarchical/upper}{0.43}
\setresult{taxonomic_alignment/plantnet/genus/independent/median}{0.01}
\setresult{taxonomic_alignment/plantnet/genus/independent/lower}{-0.14}
\setresult{taxonomic_alignment/plantnet/genus/independent/upper}{0.06}
\setresult{taxonomic_alignment/plantnet/genus/conditional/median}{-0.07}
\setresult{taxonomic_alignment/plantnet/genus/conditional/lower}{-0.21}
\setresult{taxonomic_alignment/plantnet/genus/conditional/upper}{0.09}
\setresult{taxonomic_alignment/plantnet/genus/conditional_path_sum/median}{0.02}
\setresult{taxonomic_alignment/plantnet/genus/conditional_path_sum/lower}{-0.16}
\setresult{taxonomic_alignment/plantnet/genus/conditional_path_sum/upper}{0.21}
\setresult{taxonomic_alignment/plantnet/family/hierarchical/median}{0.18}
\setresult{taxonomic_alignment/plantnet/family/hierarchical/lower}{0.10}
\setresult{taxonomic_alignment/plantnet/family/hierarchical/upper}{0.32}
\setresult{taxonomic_alignment/plantnet/family/independent/median}{0.03}
\setresult{taxonomic_alignment/plantnet/family/independent/lower}{-0.09}
\setresult{taxonomic_alignment/plantnet/family/independent/upper}{0.08}
\setresult{taxonomic_alignment/plantnet/family/conditional/median}{-0.07}
\setresult{taxonomic_alignment/plantnet/family/conditional/lower}{-0.28}
\setresult{taxonomic_alignment/plantnet/family/conditional/upper}{0.04}
\setresult{taxonomic_alignment/plantnet/family/conditional_path_sum/median}{0.22}
\setresult{taxonomic_alignment/plantnet/family/conditional_path_sum/lower}{0.04}
\setresult{taxonomic_alignment/plantnet/family/conditional_path_sum/upper}{0.49}
\setresult{taxonomic_alignment/plantnet/order/hierarchical/median}{0.17}
\setresult{taxonomic_alignment/plantnet/order/hierarchical/lower}{0.12}
\setresult{taxonomic_alignment/plantnet/order/hierarchical/upper}{0.26}
\setresult{taxonomic_alignment/plantnet/order/independent/median}{0.02}
\setresult{taxonomic_alignment/plantnet/order/independent/lower}{-0.07}
\setresult{taxonomic_alignment/plantnet/order/independent/upper}{0.08}
\setresult{taxonomic_alignment/plantnet/order/conditional/median}{-0.07}
\setresult{taxonomic_alignment/plantnet/order/conditional/lower}{-0.24}
\setresult{taxonomic_alignment/plantnet/order/conditional/upper}{0.03}
\setresult{taxonomic_alignment/plantnet/order/conditional_path_sum/median}{0.24}
\setresult{taxonomic_alignment/plantnet/order/conditional_path_sum/lower}{0.05}
\setresult{taxonomic_alignment/plantnet/order/conditional_path_sum/upper}{0.43}
\setresult{taxonomic_alignment/plantnet/class/hierarchical/median}{0.02}
\setresult{taxonomic_alignment/plantnet/class/hierarchical/lower}{-0.02}
\setresult{taxonomic_alignment/plantnet/class/hierarchical/upper}{0.16}
\setresult{taxonomic_alignment/plantnet/class/independent/median}{-0.00}
\setresult{taxonomic_alignment/plantnet/class/independent/lower}{-0.03}
\setresult{taxonomic_alignment/plantnet/class/independent/upper}{0.03}
\setresult{taxonomic_alignment/plantnet/class/conditional/median}{-0.00}
\setresult{taxonomic_alignment/plantnet/class/conditional/lower}{-0.04}
\setresult{taxonomic_alignment/plantnet/class/conditional/upper}{0.03}
\setresult{taxonomic_alignment/plantnet/class/conditional_path_sum/median}{0.39}
\setresult{taxonomic_alignment/plantnet/class/conditional_path_sum/lower}{0.15}
\setresult{taxonomic_alignment/plantnet/class/conditional_path_sum/upper}{0.52}
\setresult{taxonomic_alignment/global_lepi/genus/hierarchical/median}{0.18}
\setresult{taxonomic_alignment/global_lepi/genus/hierarchical/lower}{0.10}
\setresult{taxonomic_alignment/global_lepi/genus/hierarchical/upper}{0.38}
\setresult{taxonomic_alignment/global_lepi/genus/independent/median}{0.07}
\setresult{taxonomic_alignment/global_lepi/genus/independent/lower}{0.00}
\setresult{taxonomic_alignment/global_lepi/genus/independent/upper}{0.23}
\setresult{taxonomic_alignment/global_lepi/genus/conditional/median}{0.09}
\setresult{taxonomic_alignment/global_lepi/genus/conditional/lower}{-0.07}
\setresult{taxonomic_alignment/global_lepi/genus/conditional/upper}{0.34}
\setresult{taxonomic_alignment/global_lepi/genus/conditional_path_sum/median}{0.23}
\setresult{taxonomic_alignment/global_lepi/genus/conditional_path_sum/lower}{-0.37}
\setresult{taxonomic_alignment/global_lepi/genus/conditional_path_sum/upper}{0.37}
\setresult{taxonomic_alignment/global_lepi/family/hierarchical/median}{0.09}
\setresult{taxonomic_alignment/global_lepi/family/hierarchical/lower}{0.03}
\setresult{taxonomic_alignment/global_lepi/family/hierarchical/upper}{0.18}
\setresult{taxonomic_alignment/global_lepi/family/independent/median}{0.07}
\setresult{taxonomic_alignment/global_lepi/family/independent/lower}{0.01}
\setresult{taxonomic_alignment/global_lepi/family/independent/upper}{0.14}
\setresult{taxonomic_alignment/global_lepi/family/conditional/median}{0.22}
\setresult{taxonomic_alignment/global_lepi/family/conditional/lower}{0.08}
\setresult{taxonomic_alignment/global_lepi/family/conditional/upper}{0.32}
\setresult{taxonomic_alignment/global_lepi/family/conditional_path_sum/median}{0.32}
\setresult{taxonomic_alignment/global_lepi/family/conditional_path_sum/lower}{-0.15}
\setresult{taxonomic_alignment/global_lepi/family/conditional_path_sum/upper}{0.47}

}{}

% Then:  \result[experiment_1, accuracy, flat]   \resulttable{name}   \autoresulttable{name}   \resultgraphic[width=\linewidth]{name}
\usepackage{pgfkeys}
\usepackage{graphicx}
\ExplSyntaxOn
\clist_new:N \l__results_keys_clist
\tl_new:N \l__results_path_tl

% \setresult{a/b/c}{value}: written by the analysis code (save_results()).
\NewDocumentCommand{\setresult}{mm}{ \pgfkeys{/results/#1/.initial={#2}} }

% \result[a, b, c]: the value stored under a/b/c. A hyphen in a key may be written for an underscore,
% e.g. \result[embedding-caption, global-lepi, text].
\NewDocumentCommand{\result}{r[]}
  {
    \clist_set:Nn \l__results_keys_clist {#1}
    \tl_set:Nx \l__results_path_tl
      { /results/ \clist_use:Nn \l__results_keys_clist { / } }
    % hyphens in the keys stand for underscores (avoids the editor's "_ must be inside math mode" warning)
    % (any `_` will do: only the character code counts in a control-sequence name)
    \tl_replace_all:Nnn \l__results_path_tl { - } { _ }
    \pgfkeysifdefined { \l__results_path_tl }
      { \exp_args:NV \pgfkeysvalueof \l__results_path_tl }
      {
        \msg_warning:nnx { results } { missing } { \l__results_path_tl }
        \textbf{??}
      }
  }
\msg_new:nnn { results } { missing } { Result~'#1'~is~not~defined. }
\ExplSyntaxOff

% \resulttable{name}: the tabular of results/tables/name.tex, to put in a table environment with its own
% placement, \caption and \label. \resultcaption{name}: the generated caption text of that table (to use or adapt).
\newcommand{\resulttable}[1]{\input{results/tables/#1}}
\newcommand{\resultcaption}[1]{\result[caption_#1, text]}
% \autoresulttable{name}[placement][label]: the whole table, i.e. a table environment with \centering, the
% tabular, the generated caption and \label{tab:name}. Optional: placement (default: none) and label
% (default: tab:name), e.g. \autoresulttable{calibration_verdicts}[htb][tab:calibration].
% Note: a "[" directly after the command is read as an option.
\NewDocumentCommand{\autoresulttable}{m o o}
  {
    \IfValueTF{#2}{\begin{table}[#2]}{\begin{table}}
      \centering
      \resulttable{#1}
      \caption{\resultcaption{#1}}
      \IfValueTF{#3}{\label{#3}}{\label{tab:#1}}
    \end{table}
  }
% \resultsection{name}: a section from results/sections/name.tex.
\newcommand{\resultsection}[1]{\input{results/sections/#1}}
% \resultgraphic[options]{name}: a figure from results/figures/name.pdf.
\newcommand{\resultgraphic}[2][]{\includegraphics[#1]{results/figures/#2}}

\IfFileExists{results/values/index.tex}{% Generated; do not edit.
\setresult{calibration_error/global_lepi/higher/flat/error}{4.28}
\setresult{calibration_error/global_lepi/higher/hierarchical/error}{1.63}
\setresult{calibration_error/global_lepi/higher/hierarchical/difference}{$-2.65$ ($-2.70$ to $-2.59$)}
\setresult{calibration_error/global_lepi/higher/hierarchical/verdict}{better}
\setresult{calibration_error/global_lepi/higher/hierarchical/backbone_verdicts}{3 better}
\setresult{calibration_error/global_lepi/higher/conditional/error}{2.28}
\setresult{calibration_error/global_lepi/higher/conditional/difference}{$-2.00$ ($-2.06$ to $-1.93$)}
\setresult{calibration_error/global_lepi/higher/conditional/verdict}{better}
\setresult{calibration_error/global_lepi/higher/conditional/backbone_verdicts}{3 better}
\setresult{calibration_error/global_lepi/higher/independent/error}{2.97}
\setresult{calibration_error/global_lepi/higher/independent/difference}{$-1.31$ ($-1.36$ to $-1.25$)}
\setresult{calibration_error/global_lepi/higher/independent/verdict}{better}
\setresult{calibration_error/global_lepi/higher/independent/backbone_verdicts}{3 better}
\setresult{calibration_error/global_lepi/species/flat/error}{2.39}
\setresult{calibration_error/global_lepi/species/hierarchical/error}{1.91}
\setresult{calibration_error/global_lepi/species/hierarchical/difference}{$-0.48$ ($-0.55$ to $-0.40$)}
\setresult{calibration_error/global_lepi/species/hierarchical/verdict}{better}
\setresult{calibration_error/global_lepi/species/hierarchical/backbone_verdicts}{2 better and 1 worse}
\setresult{calibration_error/global_lepi/species/conditional/error}{5.09}
\setresult{calibration_error/global_lepi/species/conditional/difference}{$2.71$ ($2.63$ to $2.79$)}
\setresult{calibration_error/global_lepi/species/conditional/verdict}{worse}
\setresult{calibration_error/global_lepi/species/conditional/backbone_verdicts}{1 better and 2 worse}
\setresult{calibration_error/global_lepi/species/independent/error}{2.45}
\setresult{calibration_error/global_lepi/species/independent/difference}{$0.06$ ($-0.01$ to $0.14$)}
\setresult{calibration_error/global_lepi/species/independent/verdict}{inconclusive}
\setresult{calibration_error/global_lepi/species/independent/backbone_verdicts}{1 better, 1 worse and 1 inconclusive}
\setresult{calibration_error/global_lepi/backbones}{3}
\setresult{calibration_error/plantnet/higher/flat/error}{5.24}
\setresult{calibration_error/plantnet/higher/hierarchical/error}{2.08}
\setresult{calibration_error/plantnet/higher/hierarchical/difference}{$-3.16$ ($-3.62$ to $-2.70$)}
\setresult{calibration_error/plantnet/higher/hierarchical/verdict}{better}
\setresult{calibration_error/plantnet/higher/hierarchical/backbone_verdicts}{2 better and 1 inconclusive}
\setresult{calibration_error/plantnet/higher/conditional/error}{3.09}
\setresult{calibration_error/plantnet/higher/conditional/difference}{$-2.15$ ($-2.57$ to $-1.72$)}
\setresult{calibration_error/plantnet/higher/conditional/verdict}{better}
\setresult{calibration_error/plantnet/higher/conditional/backbone_verdicts}{2 better and 1 worse}
\setresult{calibration_error/plantnet/higher/independent/error}{2.61}
\setresult{calibration_error/plantnet/higher/independent/difference}{$-2.63$ ($-3.06$ to $-2.20$)}
\setresult{calibration_error/plantnet/higher/independent/verdict}{better}
\setresult{calibration_error/plantnet/higher/independent/backbone_verdicts}{2 better and 1 inconclusive}
\setresult{calibration_error/plantnet/species/flat/error}{3.99}
\setresult{calibration_error/plantnet/species/hierarchical/error}{3.36}
\setresult{calibration_error/plantnet/species/hierarchical/difference}{$-0.63$ ($-1.03$ to $-0.23$)}
\setresult{calibration_error/plantnet/species/hierarchical/verdict}{better}
\setresult{calibration_error/plantnet/species/hierarchical/backbone_verdicts}{2 better and 1 inconclusive}
\setresult{calibration_error/plantnet/species/conditional/error}{11.85}
\setresult{calibration_error/plantnet/species/conditional/difference}{$7.86$ ($7.44$ to $8.28$)}
\setresult{calibration_error/plantnet/species/conditional/verdict}{worse}
\setresult{calibration_error/plantnet/species/conditional/backbone_verdicts}{3 worse}
\setresult{calibration_error/plantnet/species/independent/error}{5.33}
\setresult{calibration_error/plantnet/species/independent/difference}{$1.34$ ($0.94$ to $1.74$)}
\setresult{calibration_error/plantnet/species/independent/verdict}{worse}
\setresult{calibration_error/plantnet/species/independent/backbone_verdicts}{2 worse and 1 inconclusive}
\setresult{calibration_error/plantnet/backbones}{3}
\setresult{calibration_error/overall/hierarchical}{9 better, 1 worse and 2 inconclusive}
\setresult{calibration_error/overall/conditional}{6 better and 6 worse}
\setresult{calibration_error/overall/independent}{6 better, 3 worse and 3 inconclusive}

% Generated; do not edit.
\setresult{caption_appendix_calibration_global_lepi/text}{Calibration error (\%) of the heads in the Global Lepidoptera dataset, per rank and backbone. Calibration error is the square root of the debiased mean squared accuracy--confidence gap across confidence bins. Better/worse: the pointwise 95\% confidence interval for the paired difference in squared calibration error lies entirely below/above 0; inconclusive: it includes 0. Bold: better calibrated than the flat head; italics: worse calibrated; plain entries (not including the flat column): inconclusive. Nonpositive estimates of the squared calibration error are displayed as zero; this does not establish exact calibration.}

% Generated; do not edit.
\setresult{caption_appendix_calibration_plantnet/text}{Calibration error (\%) of the heads in the Pl@ntNet300K dataset, per rank and backbone. Calibration error is the square root of the debiased mean squared accuracy--confidence gap across confidence bins. Better/worse: the pointwise 95\% confidence interval for the paired difference in squared calibration error lies entirely below/above 0; inconclusive: it includes 0. Bold: better calibrated than the flat head; italics: worse calibrated; plain entries (not including the flat column): inconclusive. Nonpositive estimates of the squared calibration error are displayed as zero; this does not establish exact calibration.}

% Generated; do not edit.
\setresult{caption_appendix_comparison_flemming_calibrated/text}{Test metrics for Expert labeled AMI (calibrated thresholds), in percent: estimate on the reporting partition (90\% of the test set) $\pm$ standard deviation over 1000 bootstrap resamples. Bold: the head is a plausible winner, i.e. its share of resamples in which it is best plausibly reaches chance ($1/K$ for $K$ heads; posterior probability above 5\%). Underlined: clear winner, i.e. every other head is below chance with posterior probability above 95\%. The probabilities concern the resampling shares given the fixed test data and predictions. Wins are decided on metrics rounded to 0.1 percentage points. Compared within each backbone and rank.}

% Generated; do not edit.
\setresult{caption_appendix_comparison_flemming_zero/text}{Test metrics for Expert labeled AMI (zero thresholds), in percent: estimate on the reporting partition (90\% of the test set) $\pm$ standard deviation over 1000 bootstrap resamples. Bold: the head is a plausible winner, i.e. its share of resamples in which it is best plausibly reaches chance ($1/K$ for $K$ heads; posterior probability above 5\%). Underlined: clear winner, i.e. every other head is below chance with posterior probability above 95\%. The probabilities concern the resampling shares given the fixed test data and predictions. Wins are decided on metrics rounded to 0.1 percentage points. Compared within each backbone and rank.}

% Generated; do not edit.
\setresult{caption_appendix_comparison_global_lepi_calibrated/text}{Test metrics for Global Lepidoptera (calibrated thresholds), in percent: estimate on the reporting partition (90\% of the test set) $\pm$ standard deviation over 200 bootstrap resamples. Bold: the head is a plausible winner, i.e. its share of resamples in which it is best plausibly reaches chance ($1/K$ for $K$ heads; posterior probability above 5\%). Underlined: clear winner, i.e. every other head is below chance with posterior probability above 95\%. The probabilities concern the resampling shares given the fixed test data and predictions. Wins are decided on metrics rounded to 0.1 percentage points. Compared within each backbone and rank.}

% Generated; do not edit.
\setresult{caption_appendix_comparison_global_lepi_zero/text}{Test metrics for Global Lepidoptera (zero thresholds), in percent: estimate on the reporting partition (90\% of the test set) $\pm$ standard deviation over 200 bootstrap resamples. Bold: the head is a plausible winner, i.e. its share of resamples in which it is best plausibly reaches chance ($1/K$ for $K$ heads; posterior probability above 5\%). Underlined: clear winner, i.e. every other head is below chance with posterior probability above 95\%. The probabilities concern the resampling shares given the fixed test data and predictions. Wins are decided on metrics rounded to 0.1 percentage points. Compared within each backbone and rank.}

% Generated; do not edit.
\setresult{caption_appendix_comparison_plantnet_calibrated/text}{Test metrics for Pl@ntNet300K (calibrated thresholds), in percent: estimate on the reporting partition (90\% of the test set) $\pm$ standard deviation over 1000 bootstrap resamples. Bold: the head is a plausible winner, i.e. its share of resamples in which it is best plausibly reaches chance ($1/K$ for $K$ heads; posterior probability above 5\%). Underlined: clear winner, i.e. every other head is below chance with posterior probability above 95\%. The probabilities concern the resampling shares given the fixed test data and predictions. Wins are decided on metrics rounded to 0.1 percentage points. Compared within each backbone and rank.}

% Generated; do not edit.
\setresult{caption_appendix_comparison_plantnet_zero/text}{Test metrics for Pl@ntNet300K (zero thresholds), in percent: estimate on the reporting partition (90\% of the test set) $\pm$ standard deviation over 1000 bootstrap resamples. Bold: the head is a plausible winner, i.e. its share of resamples in which it is best plausibly reaches chance ($1/K$ for $K$ heads; posterior probability above 5\%). Underlined: clear winner, i.e. every other head is below chance with posterior probability above 95\%. The probabilities concern the resampling shares given the fixed test data and predictions. Wins are decided on metrics rounded to 0.1 percentage points. Compared within each backbone and rank.}

% Generated; do not edit.
\setresult{caption_appendix_win_loss/text}{Share (\%) of comparisons in which each head is a winner or a loser, and a clear winner or loser (see Methods), at the species, genus and family rank. A comparison is one dataset, backbone, setting and metric (n = 108 per rank, excluding coverage). Mean: average over the three ranks.}

% Generated; do not edit.
\setresult{caption_calibration_verdicts/text}{Calibration of the heads compared with the flat head: the number of cases (3 backbones in each of 2 datasets) in which a head is better or worse calibrated than the flat head, at the species rank and pooled over the higher ranks. Calibration error is the square root of the debiased mean squared accuracy--confidence gap across confidence bins. Better/worse: the pointwise 95\% confidence interval for the paired difference in squared calibration error lies entirely below/above 0; inconclusive: it includes 0.}

% Generated; do not edit.
\setresult{caption_main_comparison_calibrated/text}{Test metrics for EfficientNetV2-Small (calibrated thresholds), in percent: estimate on the reporting partition (90\% of the test set) $\pm$ standard deviation over 1000 (Pl@ntNet300K and Expert labeled AMI) or 200 (Global Lepidoptera) bootstrap resamples. Bold: the head is a plausible winner, i.e. its share of resamples in which it is best plausibly reaches chance ($1/K$ for $K$ heads; posterior probability above 5\%). Underlined: clear winner, i.e. every other head is below chance with posterior probability above 95\%. The probabilities concern the resampling shares given the fixed test data and predictions. Wins are decided on metrics rounded to 0.1 percentage points. Compared within each dataset and rank.}

% Generated; do not edit.
\setresult{caption_main_comparison_zero/text}{Test metrics for EfficientNetV2-Small (zero thresholds), in percent: estimate on the reporting partition (90\% of the test set) $\pm$ standard deviation over 1000 (Pl@ntNet300K and Expert labeled AMI) or 200 (Global Lepidoptera) bootstrap resamples. Bold: the head is a plausible winner, i.e. its share of resamples in which it is best plausibly reaches chance ($1/K$ for $K$ heads; posterior probability above 5\%). Underlined: clear winner, i.e. every other head is below chance with posterior probability above 95\%. The probabilities concern the resampling shares given the fixed test data and predictions. Wins are decided on metrics rounded to 0.1 percentage points. Compared within each dataset and rank.}

% Generated; do not edit.
\setresult{embedding_caption/plantnet/species}{34}
\setresult{embedding_caption/plantnet/rare}{11}
\setresult{embedding_caption/plantnet/without_photo}{1}
\setresult{embedding_caption/plantnet/cap}{30}
\setresult{embedding_caption/plantnet/min_images}{4}
\setresult{embedding_caption/plantnet/text}{Embedding structure of the 34 species of the family \textit{Ranunculaceae} in the Pl@ntNet300K dataset: up to 30 test images per species and the species' class means, at the angle to the family's spherical median (rings every 15 degrees). The radial axis of each panel runs from 50 to 75 degrees (Flat), 45 to 75 degrees (Hierarchical), 45 to 70 degrees (Independent) and 5 to 45 degrees (Conditional); images beyond the outer limit are placed on it. Hollow markers (faded in the photo figures) show the 11 species with fewer than 4 images. No photo is available for 1 species.}
\setresult{embedding_caption/global_lepi/species}{68}
\setresult{embedding_caption/global_lepi/rare}{0}
\setresult{embedding_caption/global_lepi/without_photo}{2}
\setresult{embedding_caption/global_lepi/cap}{30}
\setresult{embedding_caption/global_lepi/min_images}{4}
\setresult{embedding_caption/global_lepi/text}{Embedding structure of the 68 species of the family \textit{Pterophoridae} in the Global Lepidoptera dataset: up to 30 test images per species and the species' class means, at the angle to the family's spherical median (rings every 15 degrees). The radial axis of each panel runs from 30 to 90 degrees (Flat), 30 to 90 degrees (Hierarchical), 35 to 80 degrees (Independent) and 30 to 75 degrees (Conditional); images beyond the outer limit are placed on it. No photo is available for 2 species.}

% Generated; do not edit.
\setresult{embedding_species_labels/plantnet}{Species labels in the Pl@ntNet300K photo figures, with the full names and GBIF taxon keys of the species: \textit{A. aestivalis} = \textit{Adonis aestivalis} (5371707); \textit{A. annua} = \textit{Adonis annua} (5371709); \textit{A. flammea} = \textit{Adonis flammea} (8246651); \textit{A. microcarpa} = \textit{Adonis microcarpa} (5616463); \textit{A. pyrenaica} = \textit{Adonis pyrenaica} (8158904); \textit{A. vernalis} = \textit{Adonis vernalis} (5371711); \textit{A. canadense} = \textit{Anemonastrum canadense} (9419395); \textit{A. narcissiflorum} = \textit{Anemonastrum narcissiflorum} (8870733); \textit{A. apennina} = \textit{Anemone apennina} (3927746); \textit{A. baldensis} = \textit{Anemone baldensis} (3927618); \textit{A. blanda} = \textit{Anemone blanda} (3033242); \textit{A. coronaria} = \textit{Anemone coronaria} (3033261); \textit{A. hortensis} = \textit{Anemone hortensis} (3925918); \textit{A. multifida} = \textit{Anemone multifida} (3033252); \textit{A. nemorosa} = \textit{Anemone nemorosa} (3033263); \textit{A. palmata} = \textit{Anemone palmata} (3924291); \textit{A. ranunculoides} = \textit{Anemone ranunculoides} (3033266); \textit{A. sylvestris} = \textit{Anemone sylvestris} (3922888); \textit{A. trifolia} = \textit{Anemone trifolia} (3033231); \textit{A. virginiana} = \textit{Anemone virginiana} (3033255); \textit{E. hupehensis} = \textit{Eriocapitella hupehensis} (9677963); \textit{E. hybrida} = \textit{Eriocapitella hybrida} (9673106); \textit{E. vitifolia} = \textit{Eriocapitella vitifolia} (3925297); \textit{H. nobilis} = \textit{Hepatica nobilis} (5371699); \textit{M. minimus} = \textit{Myosurus minimus} (3033605); \textit{P. alpina} = \textit{Pulsatilla alpina} (7278071); \textit{P. halleri} = \textit{Pulsatilla halleri} (3929423); \textit{P. montana} = \textit{Pulsatilla montana} (7277925); \textit{P. patens} = \textit{Pulsatilla patens} (3033623); \textit{P. pratensis} = \textit{Pulsatilla pratensis} (7277878); \textit{P. rubra} = \textit{Pulsatilla rubra} (7361164); \textit{P. vernalis} = \textit{Pulsatilla vernalis} (3928615); \textit{P. vulgaris} = \textit{Pulsatilla vulgaris} (3033621).}
\setresult{embedding_species_labels/global_lepi}{Species labels in the Global Lepidoptera photo figures, with the full names and GBIF taxon keys of the species: \textit{A. ambrosiae} = \textit{Adaina ambrosiae} (1858706); \textit{A. microdactyla} = \textit{Adaina microdactyla} (1858698); \textit{A. montanus} = \textit{Adaina montanus} (9689122); \textit{A. acanthadactyla} = \textit{Amblyptilia acanthadactyla} (1859285); \textit{A. falcatalis} = \textit{Amblyptilia falcatalis} (9590272); \textit{A. pica} = \textit{Amblyptilia pica} (1859287); \textit{A. punctidactyla} = \textit{Amblyptilia punctidactyla} (1859305); \textit{A. repletalis} = \textit{Amblyptilia repletalis} (9801765); \textit{A. marmarodactyla} = \textit{Anstenoptilia marmarodactyla} (1858615); \textit{C. rhododactyla} = \textit{Cnaemidophorus rhododactyla} (1858637); \textit{C. aglaodesma} = \textit{Cosmoclostis aglaodesma} (1858973); \textit{C. distans} = \textit{Crombrugghia distans} (1860063); \textit{C. tristis} = \textit{Crombrugghia tristis} (1860076); \textit{D. californicus} = \textit{Dejongia californicus} (1858659); \textit{D. lobidactylus} = \textit{Dejongia lobidactylus} (1858660); \textit{D. fasciola} = \textit{Diacrotricha fasciola} (1859453); \textit{E. monodactyla} = \textit{Emmelina monodactyla} (1859046); \textit{E. pumilio} = \textit{Exelastis pumilio} (6133397); \textit{G. didactyla} = \textit{Geina didactyla} (5123042); \textit{G. periscelidactylus} = \textit{Geina periscelidactylus} (5123046); \textit{G. ochrodactyla} = \textit{Gillmeria ochrodactyla} (4536483); \textit{G. pallidactyla} = \textit{Gillmeria pallidactyla} (4536472); \textit{H. balanotes} = \textit{Hellinsia balanotes} (9675743); \textit{H. chamelai} = \textit{Hellinsia chamelai} (8790054); \textit{H. elliottii} = \textit{Hellinsia elliottii} (9746673); \textit{H. homodactylus} = \textit{Hellinsia homodactylus} (9748079); \textit{H. inquinatus} = \textit{Hellinsia inquinatus} (9686431); \textit{H. osteodactylus} = \textit{Hellinsia osteodactylus} (1859480); \textit{H. tephradactyla} = \textit{Hellinsia tephradactyla} (1859488); \textit{L. pusillidactylus} = \textit{Lantanophaga pusillidactylus} (1858674); \textit{L. albistriolatus} = \textit{Lioptilodes albistriolatus} (9760511); \textit{M. lunaedactyla} = \textit{Marasmarcha lunaedactyla} (1859534); \textit{M. leucodactylus} = \textit{Megalorhipida leucodactylus} (1858930); \textit{M. baliodactylus} = \textit{Merrifieldia baliodactylus} (7642240); \textit{M. leucodactyla} = \textit{Merrifieldia leucodactyla} (8200050); \textit{M. indentatus} = \textit{Michaelophorus indentatus} (9700261); \textit{N. cinctipedalis} = \textit{Nippoptilia cinctipedalis} (9749806); \textit{O. eupatorii} = \textit{Oidaematophorus eupatorii} (9752950); \textit{O. lithodactyla} = \textit{Oidaematophorus lithodactyla} (1859952); \textit{O. pilosellae} = \textit{Oxyptilus pilosellae} (1858873); \textit{P. chlorosata} = \textit{Petrophora chlorosata} (1973467); \textit{P. convergata} = \textit{Petrophora convergata} (4525132); \textit{P. divisata} = \textit{Petrophora divisata} (1973458); \textit{P. narbonea} = \textit{Petrophora narbonea} (12046692); \textit{P. subaequaria} = \textit{Petrophora subaequaria} (1973466); \textit{P. calodactyla} = \textit{Platyptilia calodactyla} (1859100); \textit{P. carduidactylus} = \textit{Platyptilia carduidactylus} (1859161); \textit{P. celidotus} = \textit{Platyptilia celidotus} (9768621); \textit{P. gonodactyla} = \textit{Platyptilia gonodactyla} (1859177); \textit{P. isodactylus} = \textit{Platyptilia isodactylus} (4536292); \textit{P. belfragei} = \textit{Pselnophorus belfragei} (9777614); \textit{P. furcatalis} = \textit{Pterophorus furcatalis} (5123118); \textit{P. innotatalis} = \textit{Pterophorus innotatalis} (5123137); \textit{P. lacteipennis} = \textit{Pterophorus lacteipennis} (5123180); \textit{P. monospilalis} = \textit{Pterophorus monospilalis} (5123276); \textit{P. pentadactyla} = \textit{Pterophorus pentadactyla} (5123221); \textit{S. alternaria} = \textit{Singularia alternaria} (9645309); \textit{S. emissalis} = \textit{Sinpunctiptilia emissalis} (9759340); \textit{S. anisodactylus} = \textit{Sphenarches anisodactylus} (1859932); \textit{S. ontario} = \textit{Sphenarches ontario} (1859929); \textit{S. xerodes} = \textit{Stangeia xerodes} (1859061); \textit{S. bipunctidactyla} = \textit{Stenoptilia bipunctidactyla} (1858778); \textit{S. pterodactyla} = \textit{Stenoptilia pterodactyla} (1858775); \textit{S. zophodactylus} = \textit{Stenoptilia zophodactylus} (1858762); \textit{S. taprobanes} = \textit{Stenoptilodes taprobanes} (1860026); \textit{W. spilodactylus} = \textit{Wheeleria spilodactylus} (6097269).}

% Generated; do not edit.
\setresult{embedding_taxa/plantnet}{Taxon names in the Pl@ntNet300K figures with their GBIF taxon keys: the family \textit{Ranunculaceae} (2410) and the genera \textit{Adonis} (3033821), \textit{Anemonastrum} (7277728), \textit{Anemone} (3016227), \textit{Eriocapitella} (9777080), \textit{Hepatica} (3033655), \textit{Myosurus} (3033602) and \textit{Pulsatilla} (3033620).}
\setresult{embedding_taxa/global_lepi}{Taxon names in the Global Lepidoptera figures with their GBIF taxon keys: the family \textit{Pterophoridae} (8863) and the genera \textit{Adaina} (1858684), \textit{Amblyptilia} (1859284), \textit{Dejongia} (1858658), \textit{Geina} (9077410), \textit{Gillmeria} (4406145), \textit{Hellinsia} (1859454), \textit{Oidaematophorus} (1859936), \textit{Platyptilia} (1859082), \textit{Pterophorus} (8331787), \textit{Sphenarches} (1859925) and \textit{Stenoptilia} (1858722). The genera grouped as Other are \textit{Anstenoptilia} (1858614), \textit{Cnaemidophorus} (1858636), \textit{Cosmoclostis} (1858971), \textit{Crombrugghia} (1860060), \textit{Diacrotricha} (1859451), \textit{Emmelina} (1859036), \textit{Exelastis} (1859547), \textit{Hexadactilia} (1858716), \textit{Lantanophaga} (1858672), \textit{Lioptilodes} (1859311), \textit{Marasmarcha} (1859503), \textit{Megalorhipida} (1858921), \textit{Merrifieldia} (7641205), \textit{Michaelophorus} (4708885), \textit{Nippoptilia} (1859540), \textit{Oxyptilus} (1858866), \textit{Pselnophorus} (1858640), \textit{Singularia} (1859070), \textit{Sinpunctiptilia} (9260167), \textit{Stangeia} (1859059), \textit{Stenoptilodes} (1860010) and \textit{Wheeleria} (8744407).}

% Generated; do not edit.
\setresult{hierarchy_contrast/images}{31088}
\setresult{hierarchy_contrast/effect/architecture}{-2.35}
\setresult{hierarchy_contrast/effect/seed42}{-2.01}
\setresult{hierarchy_contrast/effect/seed43}{-2.35}
\setresult{hierarchy_contrast/effect/seed44}{-2.34}
\setresult{hierarchy_contrast/difference/estimate}{-0.12}
\setresult{hierarchy_contrast/difference/lower}{-0.54}
\setresult{hierarchy_contrast/difference/upper}{0.30}

% Generated; do not edit.
\setresult{main_comparison_win_loss/conditional/family/wins}{19.44}
\setresult{main_comparison_win_loss/conditional/family/clear_winner}{4.63}
\setresult{main_comparison_win_loss/conditional/family/winner}{14.81}
\setresult{main_comparison_win_loss/conditional/family/losses}{43.52}
\setresult{main_comparison_win_loss/conditional/family/clear_loser}{23.15}
\setresult{main_comparison_win_loss/conditional/family/loser}{20.37}
\setresult{main_comparison_win_loss/conditional/genus/wins}{21.30}
\setresult{main_comparison_win_loss/conditional/genus/clear_winner}{7.41}
\setresult{main_comparison_win_loss/conditional/genus/winner}{13.89}
\setresult{main_comparison_win_loss/conditional/genus/losses}{53.70}
\setresult{main_comparison_win_loss/conditional/genus/clear_loser}{40.74}
\setresult{main_comparison_win_loss/conditional/genus/loser}{12.96}
\setresult{main_comparison_win_loss/conditional/mean/wins}{19.14}
\setresult{main_comparison_win_loss/conditional/mean/clear_winner}{5.86}
\setresult{main_comparison_win_loss/conditional/mean/winner}{13.27}
\setresult{main_comparison_win_loss/conditional/mean/losses}{55.56}
\setresult{main_comparison_win_loss/conditional/mean/clear_loser}{41.98}
\setresult{main_comparison_win_loss/conditional/mean/loser}{13.58}
\setresult{main_comparison_win_loss/conditional/species/wins}{16.67}
\setresult{main_comparison_win_loss/conditional/species/clear_winner}{5.56}
\setresult{main_comparison_win_loss/conditional/species/winner}{11.11}
\setresult{main_comparison_win_loss/conditional/species/losses}{69.44}
\setresult{main_comparison_win_loss/conditional/species/clear_loser}{62.04}
\setresult{main_comparison_win_loss/conditional/species/loser}{7.41}
\setresult{main_comparison_win_loss/flat/family/wins}{22.22}
\setresult{main_comparison_win_loss/flat/family/clear_winner}{9.26}
\setresult{main_comparison_win_loss/flat/family/winner}{12.96}
\setresult{main_comparison_win_loss/flat/family/losses}{48.15}
\setresult{main_comparison_win_loss/flat/family/clear_loser}{23.15}
\setresult{main_comparison_win_loss/flat/family/loser}{25.00}
\setresult{main_comparison_win_loss/flat/genus/wins}{16.67}
\setresult{main_comparison_win_loss/flat/genus/clear_winner}{5.56}
\setresult{main_comparison_win_loss/flat/genus/winner}{11.11}
\setresult{main_comparison_win_loss/flat/genus/losses}{41.67}
\setresult{main_comparison_win_loss/flat/genus/clear_loser}{29.63}
\setresult{main_comparison_win_loss/flat/genus/loser}{12.04}
\setresult{main_comparison_win_loss/flat/mean/wins}{28.40}
\setresult{main_comparison_win_loss/flat/mean/clear_winner}{9.88}
\setresult{main_comparison_win_loss/flat/mean/winner}{18.52}
\setresult{main_comparison_win_loss/flat/mean/losses}{39.51}
\setresult{main_comparison_win_loss/flat/mean/clear_loser}{23.77}
\setresult{main_comparison_win_loss/flat/mean/loser}{15.74}
\setresult{main_comparison_win_loss/flat/species/wins}{46.30}
\setresult{main_comparison_win_loss/flat/species/clear_winner}{14.81}
\setresult{main_comparison_win_loss/flat/species/winner}{31.48}
\setresult{main_comparison_win_loss/flat/species/losses}{28.70}
\setresult{main_comparison_win_loss/flat/species/clear_loser}{18.52}
\setresult{main_comparison_win_loss/flat/species/loser}{10.19}
\setresult{main_comparison_win_loss/hierarchical/family/wins}{64.81}
\setresult{main_comparison_win_loss/hierarchical/family/clear_winner}{26.85}
\setresult{main_comparison_win_loss/hierarchical/family/winner}{37.96}
\setresult{main_comparison_win_loss/hierarchical/family/losses}{1.85}
\setresult{main_comparison_win_loss/hierarchical/family/clear_loser}{1.85}
\setresult{main_comparison_win_loss/hierarchical/family/loser}{0.00}
\setresult{main_comparison_win_loss/hierarchical/genus/wins}{58.33}
\setresult{main_comparison_win_loss/hierarchical/genus/clear_winner}{29.63}
\setresult{main_comparison_win_loss/hierarchical/genus/winner}{28.70}
\setresult{main_comparison_win_loss/hierarchical/genus/losses}{6.48}
\setresult{main_comparison_win_loss/hierarchical/genus/clear_loser}{2.78}
\setresult{main_comparison_win_loss/hierarchical/genus/loser}{3.70}
\setresult{main_comparison_win_loss/hierarchical/mean/wins}{54.63}
\setresult{main_comparison_win_loss/hierarchical/mean/clear_winner}{22.84}
\setresult{main_comparison_win_loss/hierarchical/mean/winner}{31.79}
\setresult{main_comparison_win_loss/hierarchical/mean/losses}{8.02}
\setresult{main_comparison_win_loss/hierarchical/mean/clear_loser}{3.09}
\setresult{main_comparison_win_loss/hierarchical/mean/loser}{4.94}
\setresult{main_comparison_win_loss/hierarchical/species/wins}{40.74}
\setresult{main_comparison_win_loss/hierarchical/species/clear_winner}{12.04}
\setresult{main_comparison_win_loss/hierarchical/species/winner}{28.70}
\setresult{main_comparison_win_loss/hierarchical/species/losses}{15.74}
\setresult{main_comparison_win_loss/hierarchical/species/clear_loser}{4.63}
\setresult{main_comparison_win_loss/hierarchical/species/loser}{11.11}
\setresult{main_comparison_win_loss/independent/family/wins}{37.96}
\setresult{main_comparison_win_loss/independent/family/clear_winner}{16.67}
\setresult{main_comparison_win_loss/independent/family/winner}{21.30}
\setresult{main_comparison_win_loss/independent/family/losses}{35.19}
\setresult{main_comparison_win_loss/independent/family/clear_loser}{25.00}
\setresult{main_comparison_win_loss/independent/family/loser}{10.19}
\setresult{main_comparison_win_loss/independent/genus/wins}{48.15}
\setresult{main_comparison_win_loss/independent/genus/clear_winner}{16.67}
\setresult{main_comparison_win_loss/independent/genus/winner}{31.48}
\setresult{main_comparison_win_loss/independent/genus/losses}{14.81}
\setresult{main_comparison_win_loss/independent/genus/clear_loser}{10.19}
\setresult{main_comparison_win_loss/independent/genus/loser}{4.63}
\setresult{main_comparison_win_loss/independent/mean/wins}{49.07}
\setresult{main_comparison_win_loss/independent/mean/clear_winner}{14.81}
\setresult{main_comparison_win_loss/independent/mean/winner}{34.26}
\setresult{main_comparison_win_loss/independent/mean/losses}{18.21}
\setresult{main_comparison_win_loss/independent/mean/clear_loser}{11.73}
\setresult{main_comparison_win_loss/independent/mean/loser}{6.48}
\setresult{main_comparison_win_loss/independent/species/wins}{61.11}
\setresult{main_comparison_win_loss/independent/species/clear_winner}{11.11}
\setresult{main_comparison_win_loss/independent/species/winner}{50.00}
\setresult{main_comparison_win_loss/independent/species/losses}{4.63}
\setresult{main_comparison_win_loss/independent/species/clear_loser}{0.00}
\setresult{main_comparison_win_loss/independent/species/loser}{4.63}

% Generated; do not edit.
\setresult{species_photo_contributors/plantnet}{Adam Golenia, Alena Shindyapkina, Antoine MARNAT, Christian Berg, Daniel Ellehammer Larsen, David Delon, Dennis Lausch, Denys Vynokurov, Flavien Saboureau, Galyna Mykytynets, Jason Grant, Jonathan, Katrin Simon, Kristof Zyskowski, Ljaž, Marco Mussita, Matt Berger, Oleksii Vasyliuk, Roger Butterfield, Roman A. Lagazyuk (Роман А. Лагазюк), Shaun Coleman, Tatsutomo Chin, Vicente Lopes, Yves Bas, carnifex, dcourtine, joan\_carles, mister\_bumble, mykhas, onotole, phiro, radoslawpuchalka and terrikoontz}
\setresult{species_photo_contributors/global_lepi}{Aalbert Rebergen, Alexis Tinker-Tsavalas, Amber M. King, Andrew Meeds, Andrzej Chruślak, Annika Lindqvist, Cathy Powers, Dan MacNeal, Dan Riley, David George, Diego González Dopico, Donald Hobern, Doreen, Evgenij Samarin, Francisco Farriols Sarabia, Fyodor Pudovikov, Jason Grant, Joey Bom, Johannes S., Kai Squires, Keith Wauthy, Ken-ichi Ueda, Laura Gaudette, Louis Imbeau, Michael Berardozzi, Nigel Voaden, Peter Crowcroft, RAP, Randy A Nonenmacher, Romilson Silva Lopes Junior, Sam Kieschnick, Saryu Mae 前 朝琉, Seig, Shannon Stoner, Stuart Fraser, Sylvain Montagner, Thomas Irvine, Thomas Koffel, Tommy Andriollo, adam1420, amzamz, bioshots\_jm, carey\_knox\_southern\_scales, kike Vergara, loarie, naturewatchwidow, pmrsco, portioid, rich\_fuller, sgalick, steve\_mcwilliam, Вячеслав Юсупов and Поляков Александр}

% Generated; do not edit.
\setresult{taxonomic_alignment/methods}{average linkage, complete linkage, PAM and HDBSCAN}
\setresult{taxonomic_alignment/subsamples}{1000}
\setresult{taxonomic_alignment/species_per_subsample}{512}
\setresult{taxonomic_alignment/plantnet/genus/hierarchical/median}{0.15}
\setresult{taxonomic_alignment/plantnet/genus/hierarchical/lower}{0.02}
\setresult{taxonomic_alignment/plantnet/genus/hierarchical/upper}{0.43}
\setresult{taxonomic_alignment/plantnet/genus/independent/median}{0.01}
\setresult{taxonomic_alignment/plantnet/genus/independent/lower}{-0.14}
\setresult{taxonomic_alignment/plantnet/genus/independent/upper}{0.06}
\setresult{taxonomic_alignment/plantnet/genus/conditional/median}{-0.07}
\setresult{taxonomic_alignment/plantnet/genus/conditional/lower}{-0.21}
\setresult{taxonomic_alignment/plantnet/genus/conditional/upper}{0.09}
\setresult{taxonomic_alignment/plantnet/genus/conditional_path_sum/median}{0.02}
\setresult{taxonomic_alignment/plantnet/genus/conditional_path_sum/lower}{-0.16}
\setresult{taxonomic_alignment/plantnet/genus/conditional_path_sum/upper}{0.21}
\setresult{taxonomic_alignment/plantnet/family/hierarchical/median}{0.18}
\setresult{taxonomic_alignment/plantnet/family/hierarchical/lower}{0.10}
\setresult{taxonomic_alignment/plantnet/family/hierarchical/upper}{0.32}
\setresult{taxonomic_alignment/plantnet/family/independent/median}{0.03}
\setresult{taxonomic_alignment/plantnet/family/independent/lower}{-0.09}
\setresult{taxonomic_alignment/plantnet/family/independent/upper}{0.08}
\setresult{taxonomic_alignment/plantnet/family/conditional/median}{-0.07}
\setresult{taxonomic_alignment/plantnet/family/conditional/lower}{-0.28}
\setresult{taxonomic_alignment/plantnet/family/conditional/upper}{0.04}
\setresult{taxonomic_alignment/plantnet/family/conditional_path_sum/median}{0.22}
\setresult{taxonomic_alignment/plantnet/family/conditional_path_sum/lower}{0.04}
\setresult{taxonomic_alignment/plantnet/family/conditional_path_sum/upper}{0.49}
\setresult{taxonomic_alignment/plantnet/order/hierarchical/median}{0.17}
\setresult{taxonomic_alignment/plantnet/order/hierarchical/lower}{0.12}
\setresult{taxonomic_alignment/plantnet/order/hierarchical/upper}{0.26}
\setresult{taxonomic_alignment/plantnet/order/independent/median}{0.02}
\setresult{taxonomic_alignment/plantnet/order/independent/lower}{-0.07}
\setresult{taxonomic_alignment/plantnet/order/independent/upper}{0.08}
\setresult{taxonomic_alignment/plantnet/order/conditional/median}{-0.07}
\setresult{taxonomic_alignment/plantnet/order/conditional/lower}{-0.24}
\setresult{taxonomic_alignment/plantnet/order/conditional/upper}{0.03}
\setresult{taxonomic_alignment/plantnet/order/conditional_path_sum/median}{0.24}
\setresult{taxonomic_alignment/plantnet/order/conditional_path_sum/lower}{0.05}
\setresult{taxonomic_alignment/plantnet/order/conditional_path_sum/upper}{0.43}
\setresult{taxonomic_alignment/plantnet/class/hierarchical/median}{0.02}
\setresult{taxonomic_alignment/plantnet/class/hierarchical/lower}{-0.02}
\setresult{taxonomic_alignment/plantnet/class/hierarchical/upper}{0.16}
\setresult{taxonomic_alignment/plantnet/class/independent/median}{-0.00}
\setresult{taxonomic_alignment/plantnet/class/independent/lower}{-0.03}
\setresult{taxonomic_alignment/plantnet/class/independent/upper}{0.03}
\setresult{taxonomic_alignment/plantnet/class/conditional/median}{-0.00}
\setresult{taxonomic_alignment/plantnet/class/conditional/lower}{-0.04}
\setresult{taxonomic_alignment/plantnet/class/conditional/upper}{0.03}
\setresult{taxonomic_alignment/plantnet/class/conditional_path_sum/median}{0.39}
\setresult{taxonomic_alignment/plantnet/class/conditional_path_sum/lower}{0.15}
\setresult{taxonomic_alignment/plantnet/class/conditional_path_sum/upper}{0.52}
\setresult{taxonomic_alignment/global_lepi/genus/hierarchical/median}{0.18}
\setresult{taxonomic_alignment/global_lepi/genus/hierarchical/lower}{0.10}
\setresult{taxonomic_alignment/global_lepi/genus/hierarchical/upper}{0.38}
\setresult{taxonomic_alignment/global_lepi/genus/independent/median}{0.07}
\setresult{taxonomic_alignment/global_lepi/genus/independent/lower}{0.00}
\setresult{taxonomic_alignment/global_lepi/genus/independent/upper}{0.23}
\setresult{taxonomic_alignment/global_lepi/genus/conditional/median}{0.09}
\setresult{taxonomic_alignment/global_lepi/genus/conditional/lower}{-0.07}
\setresult{taxonomic_alignment/global_lepi/genus/conditional/upper}{0.34}
\setresult{taxonomic_alignment/global_lepi/genus/conditional_path_sum/median}{0.23}
\setresult{taxonomic_alignment/global_lepi/genus/conditional_path_sum/lower}{-0.37}
\setresult{taxonomic_alignment/global_lepi/genus/conditional_path_sum/upper}{0.37}
\setresult{taxonomic_alignment/global_lepi/family/hierarchical/median}{0.09}
\setresult{taxonomic_alignment/global_lepi/family/hierarchical/lower}{0.03}
\setresult{taxonomic_alignment/global_lepi/family/hierarchical/upper}{0.18}
\setresult{taxonomic_alignment/global_lepi/family/independent/median}{0.07}
\setresult{taxonomic_alignment/global_lepi/family/independent/lower}{0.01}
\setresult{taxonomic_alignment/global_lepi/family/independent/upper}{0.14}
\setresult{taxonomic_alignment/global_lepi/family/conditional/median}{0.22}
\setresult{taxonomic_alignment/global_lepi/family/conditional/lower}{0.08}
\setresult{taxonomic_alignment/global_lepi/family/conditional/upper}{0.32}
\setresult{taxonomic_alignment/global_lepi/family/conditional_path_sum/median}{0.32}
\setresult{taxonomic_alignment/global_lepi/family/conditional_path_sum/lower}{-0.15}
\setresult{taxonomic_alignment/global_lepi/family/conditional_path_sum/upper}{0.47}

}{}
